\section{4. Mô hình và giải pháp đề xuất}

\begin{frame}[plain]
\centering
\vfill
{\usebeamerfont{title}\usebeamercolor[fg]{title}\LARGE\textbf{4. Mô hình và giải pháp đề xuất}}
\vfill
\end{frame}

% Source slide 24
\begin{frame}[t]{4.1. Tổng quan giải pháp}
\small
\begin{block}{Lớp hỗ trợ quyết định trước Slurm}
\centering
\textbf{Slurm script + \texttt{\#INTENT}}
$\rightarrow$ Intent interpretation
$\rightarrow$ Structured intent
\end{block}
\vspace{0.3em}
\begin{columns}[T,onlytextwidth]
\begin{column}{0.31\textwidth}
\begin{exampleblock}{1. Context}
Quota, queue state và workload evidence
\end{exampleblock}
\end{column}
\begin{column}{0.34\textwidth}
\begin{exampleblock}{2. Evaluate}
Candidate profiles $\rightarrow$ feasibility $\rightarrow$ prediction + uncertainty
\end{exampleblock}
\end{column}
\begin{column}{0.31\textwidth}
\begin{exampleblock}{3. Select}
Pareto/queue-aware selection và evidence
\end{exampleblock}
\end{column}
\end{columns}
\vspace{0.5em}
\begin{alertblock}{Action cuối cùng}
\centering
\texttt{commit} \quad | \quad \texttt{keep\_original} \quad | \quad \texttt{schedule\_probe}
\end{alertblock}
\centering Modified Slurm script $\rightarrow$ Slurm scheduler không thay đổi.
\end{frame}

% Source slide 25
\begin{frame}[t,shrink=20]{4.2. Mô hình đầu vào và đầu ra}
\small
Đầu vào của hệ thống:\par
\[
X_j=(S_j,U_j,Q_u,H_t,E_j)
\]
Trong đó:\par
\begin{itemize}
\item \(S_j\): Slurm script ban đầu;
\item \(U_j\): user intent;
\item \(Q_u\): quota còn lại;
\item \(H_t\): trạng thái cluster và queue;
\item \(E_j\): profiling và operational evidence.
\end{itemize}
Đầu ra:\par
\[
D_j=(action,profile,predictions,uncertainty,evidence)
\]
\end{frame}

% Source slide 26
\begin{frame}[t,shrink=20]{4.3. Mô hình ý định người dùng}
\small
Intent được biểu diễn bởi:\par
\[
I_j=(G_j,C_j,F_j,P_j,U_j^{unc})
\]
Trong đó:\par
\begin{itemize}
\item \(G_j\): mục tiêu chính;
\item \(C_j\): ràng buộc cứng;
\item \(F_j\): các thay đổi tài nguyên được user cho phép;
\item \(P_j\): các ưu tiên mềm;
\item \(U_j^{unc}\): thông tin thiếu, mơ hồ hoặc xung đột.
\end{itemize}
Mô hình này giúp tách rõ constraint, preference và flexibility.\par
\end{frame}

% Source slide 27
\begin{frame}[t,shrink=20,fragile]{4.3. Mô hình ý định người dùng}
\small
Ví dụ intent:\par
\begin{quote}
\emph{Cần hoàn thành trước 18 giờ. Có thể dùng ít GPU hơn nếu vẫn kịp deadline. Ưu tiên tiết kiệm quota.}
\end{quote}
Structured intent:\par
\begin{block}{Minh họa}
\begin{verbatim}
{
  "goal": {
    "primary": "meet_deadline",
    "secondary": "reduce_quota"
  },
  "hard_constraints": {
    "deadline": "18:00"
  },
  "authorized_flexibility": {
    "gpu_reduction_allowed": true
  },
  "preferences": [
    "lower_gpu_seconds"
  ]
}
\end{verbatim}
\end{block}
\end{frame}

% Source slide 28
\begin{frame}[t,shrink=20]{4.4. LLM và bộ kiểm tra xác định}
\small
Mô hình ngôn ngữ được sử dụng để chuyển đổi yêu cầu tự nhiên thành intent có cấu trúc.\par
Sau đó, bộ kiểm tra xác định được dùng để:\par
\begin{itemize}
\item kiểm tra schema;
\item phát hiện field không hợp lệ;
\item phát hiện conflict;
\item loại bỏ flexibility chưa được user cho phép;
\item phát hiện unsupported intent;
\item ngăn commit quá sớm.
\end{itemize}
LLM không trực tiếp quyết định profile cuối cùng.\par
\end{frame}

% Source slide 29
\begin{frame}[t,shrink=20,fragile]{4.5. Sinh candidate profiles}
\small
Candidate profiles được sinh từ các ràng buộc:\par
\[
\mathcal{P}_j^{cand}=
\{p:
C_{cluster}(p)
\land C_{workload}(p)
\land C_{intent}(p)
\land C_{quota}(p)\}
\]
Các bước chính:\par
\begin{block}{Minh họa}
\begin{verbatim}
All combinations
→ cluster-valid
→ workload-valid
→ intent-valid
→ quota-valid
→ candidate profiles
\end{verbatim}
\end{block}
Đây là constrained search, không phải free-form generation.\par
\end{frame}

% Source slide 30
\begin{frame}[t]{4.6. Feasibility và prediction uncertainty}
\small
\begin{columns}[T,onlytextwidth]
\begin{column}{0.49\textwidth}
\begin{block}{Feasibility kỹ thuật}
Feasibility có thể không liên tục theo số GPU, CPU hoặc memory. Ở biên 4 GPU:\par
\begin{itemize}
\item một số workload chạy thành công;
\item một số workload fail do tham số không tương thích trước compute.
\end{itemize}
\textbf{Cần kết hợp:}\par
\begin{itemize}
\item workload rule và preflight;
\item measured evidence và prediction.
\end{itemize}
\end{block}
\end{column}
\begin{column}{0.49\textwidth}
\begin{block}{Prediction và uncertainty}
Mỗi candidate profile cần được đánh giá bởi:\par
\begin{itemize}
\item runtime prediction;
\item memory feasibility;
\item interval hoặc upper bound;
\item relative uncertainty width.
\end{itemize}
\textbf{Mô hình khảo sát:}\par
\begin{itemize}
\item factorized và monolithic predictor;
\item kNN và conformal upper bound.
\end{itemize}
\end{block}
\end{column}
\end{columns}
\end{frame}

% Source slide 32
\begin{frame}[t,shrink=20]{4.8. Static selection và queue-aware selection}
\small
Static selection xét:\par
\begin{itemize}
\item runtime;
\item GPU-seconds;
\item quota;
\item prediction risk.
\end{itemize}
Queue-aware selection bổ sung:\par
\begin{itemize}
\item current occupancy;
\item estimated waiting time;
\item GPU fragmentation;
\item pressure state.
\end{itemize}
\[
Turnaround(p)=Waiting(p,H_t)+Runtime(p)
\]
Khi queue bị pressure, profile có runtime dài hơn nhưng wait time thấp hơn có thể tạo turnaround tốt hơn.\par
\end{frame}

% Source slide 33
\begin{frame}[t,shrink=20]{4.9. Cơ chế quyết định cuối cùng}
\small
Hệ thống có ba action:\par
\begin{itemize}
\item \textbf{commit:} đề xuất profile mới vì đủ evidence và confidence;
\item \textbf{keep original:} giữ request ban đầu vì không có profile tốt hơn đủ rõ;
\item \textbf{schedule probe:} cần đo thêm vì uncertainty còn cao.
\end{itemize}
Điểm quan trọng của giải pháp:\par
\begin{quote}
\emph{Hệ thống không bắt buộc phải tự động thay đổi request trong mọi trường hợp.}
\end{quote}
\end{frame}

% Source slide 34
\begin{frame}[t,shrink=20]{4.10. Khả năng kiểm chứng của quyết định}
\small
Mỗi quyết định cần lưu lại:\par
\begin{itemize}
\item intent đã parse;
\item ràng buộc được áp dụng;
\item candidate profiles;
\item profile bị loại và lý do;
\item prediction và interval;
\item queue snapshot;
\item action cuối cùng;
\item script sau khi rewrite;
\item kết quả thực thi sau submit.
\end{itemize}
\[
Intent
\rightarrow Evidence
\rightarrow Decision
\rightarrow Outcome
\]
\end{frame}
